Käesoleva lõputöö eesmärgiks on luua töötav näidislahendus, 
mis demonstreerib eDelivery dünaamilise otsimise mehanismi kasutust.
Lõppeesmärgini jõudmiseks tegi autor neljaosalise plaani:

\begin{enumerate}
  \item Saada tööle staatiline eDelivery konfiguratsioon.
  \item Konfigureerida juurdepääsupunkt ümber, selleks et võimaldada dünaamilist otsingut.
  \item Luua DNS tarkvara, mis suudab lahendada juurdepääsupunkti poolt tehtavaid päringuid SML-ile.
  \item Saada tööle SMP server, mis suudab vastata juurdepääsupunkti päringutele.
\end{enumerate}

\section{Staatiline eDelivery konfiguratsioon}
eDelivery staatilise konfiguratsiooni puhul vahetavad juurdepääsupunktid infot otse, 
ilma kolmanda osapooleta. See on tehniliselt vähem keerukas lahendus, kui dünaamiline otsing,
ning selle tõttu ka hea esimene samm lõppeesmärgini jõudmiseks.

\subsection{Juurdepääsupunktide konteinerite seadistamine}
Näidislahendus eeldab, et kaks juurdepääsupunkti suudavad omavahel infot jagada.
Sellest tulenevalt on vaja käivitada kaks instantsi juurdepääsupunkti tarkvara.

Tarkvara käivitamise lihtsustamiseks, ning üheselt igas masinas samamoodi töötava lahenduse loomiseks
kasutab näidislahendus Dockeri konteinereid. Kuna käivitada tuleb mitu erinevat konteinerit,
siis kasutab näidislahendus \textit{docker compose} mehhanismi. 
Compose võimaldab käivitada palju erinevaid konteinereid koos kindla konfiguratsiooniga, 
kasutades ainult ühte definitsioonifaili: \textit{compose.yml}

Harmony eDelivery juurdepääsupunkt on hästi konteineriseeritud ning dokumenteeritud tarkvara
mis on saadaval Docker Hub-is. Seal on olemas ka dokumentatsioon kuidas kasutada Harmony konteinereid
compose failis. Sealsete juhitse kohaselt nõuab üks Harmony instants ka MySQL andmebaasi. \cite{niis_harmony_ap_docker}
See tähendab, et kahe Harmony instantsi jooksutamine nõuab kokku nelja konteinerit.

Näidislahenduse \textit{compose.yml} fail sisaldab kahte juurdepääsupunkti
nimedega \textit{sender} ja \textit{receiver}, vastavalt nende tööpõhimõtetele, 
ning ka kahte MySQL andmebaasi analoogsete nimedega.

\subsection{Juurdepääsupunktide staatiline konfigureerimine}
Juurdepääsupunkte on võimalik seadistada kahest kohast - keskkonnamuutujad ja veebipõhine kasutajaliides.

Näidislahenduses on \textit{compose.yml} lisatud paar keskkonnamuutujat, mida Docker Hubis 
oleval näidisel ei olnud. Need lisandused võimaldavad juurdepääsupunkti õige toimimise.
keskkonnamuutujatest kõige olulisemad on LB\_TLS\_TERMINATION ja SECURITY\_*\_PASSWORD.

\subsubsection{Keskkonnamuutujad}
LB\_TLS\_TERMINATION muudab juurdepääsupunkti kättesaadavaks pordil 80 (http), mis muudab
digitaalsete sertifikaatide haldamise lihtsamaks, kuna ei pea arvestama TLS sertifikaatidega.

SECURITY\_*\_PASSWORD võimaldab kasutajal, läbi veebiliidese, laadida alla, 
sirvida ning muuta juurdepääsupunkti privaatseid võtmeid ning avalikke sertifikaate. 
Ligipääs avalikkele sertifikaatidele on oluline selleks, 
et neid saaks jagada juurdepääsupunktide vahel, kas staatiliselt või dünaamiliselt.